\chapter{Modelo de Domínio}

O modelo de domínio representa os objetos centrais que descrevem o ciclo de vida de uma análise de imagem. Ele foi mantido propositalmente enxuto, porém suficiente para representar a solicitação, o processamento e o resultado de uma inferência futura. O objetivo é preservar a clareza do domínio sem misturar preocupações de infraestrutura ou apresentação.

\section{AnalysisJob}

\textit{AnalysisJob} representa uma solicitação completa de análise enviada ao microserviço. O objeto reúne identificadores, metadados, referência da imagem, estado de processamento, datas de criação e conclusão, além do resultado opcional. Esse tipo é o ponto central do fluxo de negócio, pois conecta a entrada da API, a persistência temporária, o armazenamento da imagem e a futura publicação de mensagens.

\section{PredictionResult}

\textit{PredictionResult} representa o resultado consolidado da inferência. Ele contém a doença principal identificada, o grau de confiança, uma descrição textual, a versão do modelo e o instante de processamento. A presença de uma coleção de \textit{DiseasePrediction} permite que a solução evolua para cenários em que mais de uma hipótese seja relevante para análise ou auditoria.

\section{DiseasePrediction}

\textit{DiseasePrediction} representa uma predição individual retornada pelo modelo. Seu papel é fornecer granularidade para a composição do resultado final, principalmente em cenários de ranking de classes, análise top-k e futuras explicações do resultado.

\section{AnalysisStatus}

\textit{AnalysisStatus} é o enum que descreve o estado de execução de uma análise. Os estados \textit{Pending}, \textit{Queued}, \textit{Processing}, \textit{Completed} e \textit{Failed} formam a linha temporal do job, permitindo leitura clara em logs, consultas e callbacks.

\section{AnalysisMetadata}

\textit{AnalysisMetadata} é o ponto de extensibilidade para dados auxiliares enviados junto à imagem. Ele guarda informações como espécie, cultura, localização e pares chave-valor adicionais. Isso evita a necessidade de alterar o modelo principal sempre que uma nova propriedade contextual for introduzida.

\section{Relacionamentos}

O relacionamento principal do domínio parte de \textit{AnalysisJob}, que referencia \textit{AnalysisMetadata} e pode conter um \textit{PredictionResult}. O \textit{PredictionResult}, por sua vez, agrega múltiplas \textit{DiseasePrediction}. Essa composição mantém o domínio coeso e compatível com futuras evoluções do mecanismo de inferência.

\begin{figure}[H]
    \centering
    \fbox{\parbox{0.85\textwidth}{\centering Espaço reservado para o diagrama de classes do modelo de domínio.}}
    \caption{Diagrama do modelo de domínio}
    \label{fig:modelo-dominio}
\end{figure}
